
\documentclass[11pt,a4paper]{article}

\usepackage[utf8]{inputenc}
\usepackage[T1]{fontenc}
\usepackage{lmodern}
\usepackage[margin=0.6in]{geometry}
\usepackage{enumitem}
\usepackage{titlesec}
\usepackage{xcolor}
\usepackage{hyperref}

\pagestyle{empty}
\setlength{\parindent}{0pt}
\setlength{\parskip}{0pt}

\titleformat{\section}{\large\bfseries\uppercase}{\thesection}{0em}{}[\titlerule]
\titlespacing*{\section}{0pt}{10pt}{6pt}

\newcommand{\entryheader}[4]{%
    \textbf{#1} \hfill \textit{#2}\\
    \textit{#3} \hfill #4
}

\begin{document}

\begin{center}
    {\LARGE\bfseries Aswanth Bindu Ajith}\\[4pt]
    \texttt{aswanth.bindu-ajith@stud.th-deg.de} | \texttt{+49 176 4153 9896} | \texttt{Deggendorf, Germany} | \texttt{LinkedIn}
\end{center}

\section{Professional Summary}
M.Sc. student in High Performance Computing and Quantum Computing with a Mechanical Engineering foundation. Comfortable building high-performance computing, software engineering solutions — from low-level numerical code to production web features. Presented optimization findings to a 25-person seminar and received top marks. Reduced simulation runtime by 40% via vectorization on a 30,000-entity dataset. Looking for a software engineering role where I can apply my mix of engineering rigor and software development experience to solve real problems.

\section{Technical Skills}

HPC & Parallel Computing: MPI (basic), OpenMP (basic), GPU-aware computing concepts, Cluster job scheduling (SLURM basics), Performance profiling\\

Programming Languages: Python (Advanced, 4+ years), C/C++, JavaScript, Bash/Shell scripting, HTML/CSS\\

Quantum Computing: Qiskit, Cirq (basic), Quantum optimization algorithms, Variational quantum eigensolver (VQE)\\

Software Engineering: Git/GitHub, Linux (Ubuntu, WSL), Docker (basic), Agile/Scrum, Unit testing, CI/CD basics


\section{Professional Experience}

\entryheader{ Working Student Software Developer }{ Germany }{ No Border (E-commerce) }{ Feb 2025 – Present }
\begin{itemize}[leftmargin=*,nosep,topsep=2pt]
    
    \item Shipped 5+ frontend features in a 2-week sprint cycle, reducing average ticket age by 30%
    
    \item Debugged and resolved 20+ production bugs across the React/Vue stack, cutting reported error rates by 25%
    
    \item Wrote unit tests for 3 core modules, pushing test coverage from 45% to 78%
    
    \item Collaborated with 4-person cross-functional team (design, backend, QA) using GitHub Flow and daily standups
    
\end{itemize}


\section{Education}

\textbf{ M.Sc. High Performance Computing and Quantum Computing } \hfill 2024 – Present\\
Deggendorf Institute of Technology (DIT), Germany \hfill Current Average: 2.6

\\\textit{Relevant coursework:} Parallel & Distributed Computing, Numerical Optimization, Quantum Algorithms, Scientific Software Engineering, High-Dimensional Data Analysis


\\\textit{Highlights:} Optimization Methods (1.7); System Design & Applications (1.7); Software Engineering (2.3); Advanced Mathematics (2.3)


\textbf{ B.Tech Mechanical Engineering } \hfill 2019 – 2023\\
APJ Abdul Kalam Technological University, India \hfill CGPA: 6.48




\section{Projects}

\textbf{ Computing Architecture & Optimization Methods }
\begin{itemize}[leftmargin=*,nosep,topsep=2pt]
    
    \item Implemented gradient descent and genetic algorithms from scratch to optimize a 50-dimensional engineering cost function
    
    \item Benchmarked CPU vs. GPU (CUDA) performance on matrix operations, documenting 8x speedup for large matrices
    
    \item Profiled memory bottlenecks using Python's cProfile and reduced peak RAM usage by 35% through generator-based streaming
    
    \item Presented findings to a 25-person seminar, receiving top marks on the technical report
    
    \item Key metrics: 50-dimensional optimization, 8x GPU speedup, 35% RAM reduction
    
\end{itemize}

\textbf{ Large-Scale Computational Simulation & Data Analysis }
\begin{itemize}[leftmargin=*,nosep,topsep=2pt]
    
    \item Built a Python simulation engine processing 30,000+ interacting entities using vectorized NumPy operations
    
    \item Reduced runtime from 12 minutes to 7 minutes (40% speedup) by replacing nested loops with Pandas groupby and broadcasting
    
    \item Identified 3 non-obvious data correlations that shifted the team's modeling assumptions
    
    \item Packaged results into an 18-page LaTeX report with reproducible Jupyter notebooks for peer review
    
    \item Key metrics: 30,000+ entities, 40% runtime reduction, 18-page report
    
\end{itemize}

\textbf{ Thermal Analysis of Battery Modules }
\begin{itemize}[leftmargin=*,nosep,topsep=2pt]
    
    \item Modeled heat transfer across a 12-cell lithium-ion battery pack using finite difference methods in Python
    
    \item Simulated 8 operating scenarios, identifying a thermal runaway risk at >45°C under sustained 3C discharge
    
    \item Proposed a revised cooling channel geometry that improved peak temperature uniformity by 18% in simulation
    
    \item Co-authored a technical report cited by 2 subsequent student projects in the same lab
    
    \item Key metrics: 12-cell pack, 8 scenarios, 18% temperature improvement
    
\end{itemize}


\section{Research Interests}
Performance Optimization, High Performance Computing, Parallel and Distributed Systems, GPU Computing

\section{Languages}
English (Professional Working Proficiency), Malayalam (Native), German (A1 (Beginner))

\end{document}